\documentclass[12pt,a4paper]{article}
\usepackage[utf8]{vietnam}
\usepackage{amsmath,amssymb}
\usepackage{tabularx}
\usepackage{tikz}
\usepackage{xcolor}
\usepackage[top=0.6cm, bottom=0.6cm, left=0.6cm, right=0.6cm]{geometry}
\usepackage[dethi]{ex_test}

\SAOPTN{widthSA=3.2,heightSA=0.62}

\setlength{\parskip}{0pt}

% Dòng kẻ chấm mờ đều đặn chuẩn viết tay
\newcommand{\dongke}[1]{%
  \begingroup%
  \setlength{\parskip}{0pt}%
  \baselineskip=18.5pt%
  \vspace{1.0mm}%
  \foreach \i in {1,...,#1}{%
    \noindent{\color{gray!60}\dotfill}\par%
  }%
  \vspace{0.6mm}%
  \endgroup%
}

% Định nghĩa ô bài tập tự luận trong bảng
\newcommand{\oBaiTap}[3][6]{%
  \vspace{1.0mm}%
  \noindent\textbf{#2} #3\par%
  \dongke{#1}%
}

\begin{document}
\pagestyle{empty}
% =============================================================================
% ========================= TOÁN 9 - TUẦN 08 - BUỔI 02: MÃ ĐỀ 101 =============
% ========================= MẶT TRƯỚC (TRANG 1/2) ==============================
% =============================================================================

\noindent
\begin{tabularx}{\linewidth}{|X|c|}
\hline
\textbf{PHIẾU KIỂM TRA ĐỊNH KỲ TOÁN 9 - TUẦN 08 - BUỔI 02} & \textbf{Mã đề: 101 - Trang 1/2} \\
\hline
\end{tabularx}

\vspace{1.0mm}

\noindent
\begin{tabularx}{\linewidth}{|p{3.2cm}|X|}
\hline
\centering\textbf{ĐIỂM SỐ} & \centering\textbf{LỜI NHẬN XÉT CỦA THẦY CÔ} \tabularnewline
\hline
\rule{0pt}{1.8cm} & \tabularnewline
\hline
\end{tabularx}

\vspace{1.2mm}
\noindent\textbf{Câu 1 (0,25 điểm).} Giá trị của biểu thức $\sqrt{49} - \sqrt{16}$ bằng: \shortans{3}\par\vspace{0.5mm}
\noindent\textbf{Câu 2 (0,25 điểm).} Biểu thức $\sqrt{2x - 6}$ xác định khi: \shortans{x \ge 3}\par\vspace{0.5mm}
\noindent\textbf{Câu 3 (0,25 điểm).} Nghiệm của bất phương trình $3x - 9 \ge 0$ là: \shortans{x \ge 3}\par\vspace{0.5mm}
\noindent\textbf{Câu 4 (0,25 điểm).} Nghiệm của hệ $\begin{cases} 2x - y = 3 \\ x + y = 3 \end{cases}$ là $(x; y) =$ \shortans{(2; 1)}\par\vspace{0.5mm}
\noindent\textbf{Câu 5 (0,25 điểm).} Cho $\Delta ABC$ vuông tại $A$ ($AB=3, AC=4$). Khi đó $\sin C =$ \shortans{\dfrac{3}{5}}\par\vspace{0.5mm}
\noindent\textbf{Câu 6 (0,25 điểm).} $\Delta ABC$ vuông tại $A$ ($BC = 10, \widehat{B}=60^\circ$). Cạnh $AC =$ \shortans{5\sqrt{3}\text{ cm}}\par\vspace{0.5mm}
\noindent\textbf{Câu 7 (0,25 điểm).} Rút gọn biểu thức $\sqrt{(3 - \sqrt{11})^2}$ được: \shortans{\sqrt{11} - 3}\par\vspace{0.5mm}
\noindent\textbf{Câu 8 (0,25 điểm).} Giá trị của biểu thức $\sin^2 35^\circ + \sin^2 55^\circ$ bằng: \shortans{1}\par\vspace{0.5mm}
\noindent\textbf{Câu 9 (0,25 điểm).} Cây cao $8\text{ m}$ đổ bóng dài $8\text{ m}$. Góc của tia nắng là: \shortans{45^\circ}\par\vspace{0.5mm}
\noindent\textbf{Câu 10 (0,25 điểm).} Nghiệm của phương trình $\sqrt{3x + 1} = 4$ là $x =$ \shortans{5}\par\vspace{0.5mm}
\vspace{0.8mm}
\noindent\textbf{Câu 11 (2,5 điểm). Đại số tổng hợp:}
\vspace{0.6mm}

\noindent
\begin{tabularx}{\linewidth}{|X|X|}
\hline
\oBaiTap[22]{a) (1,25 điểm)}{Giải hệ phương trình:\par $\begin{cases} 2x - y = 6 \\ 3x + y = 14 \end{cases}$} & 
\oBaiTap[22]{b) (1,25 điểm)}{Rút gọn với $x > 0, x \ne 1$:\par $B = \left(\dfrac{\sqrt{x}}{\sqrt{x}-1} - \dfrac{1}{x-\sqrt{x}}\right) : \dfrac{\sqrt{x}+1}{x-1}$} \\
\hline
\end{tabularx}

\newpage
% =============================================================================
% ========================= TOÁN 9 - TUẦN 08 - BUỔI 02: MÃ ĐỀ 101 =============
% ========================= MẶT SAU (TRANG 2/2) ================================
% =============================================================================

\noindent
\begin{tabularx}{\linewidth}{|X|c|}
\hline
\textbf{TOÁN 9 - TUẦN 08 - BUỔI 02} & \textbf{Mã đề: 101 - Trang 2/2} \\
\hline
\end{tabularx}

\vspace{1.2mm}

\noindent
\begin{tabularx}{\linewidth}{|X|}
\hline
\vspace{0.5mm}
\textbf{Câu 12 (2,5 điểm).} Một đội xe dự định chở $75$ tấn hàng để ủng hộ đồng bào miền Trung bị lũ lụt. Lúc sắp khởi hành, đội nhận được ủng hộ thêm $5$ tấn hàng nữa và đồng thời được bổ sung thêm $5$ xe cùng loại, do đó mỗi xe chở ít hơn dự định $1$ tấn hàng. Hỏi ban đầu đội xe dự định sử dụng bao nhiêu chiếc xe? (Biết mỗi xe chở khối lượng hàng như nhau).\par
\vspace{0.4mm}
\noindent\textit{Bài giải:}\dongke{16}
\\ \hline
\end{tabularx}

\vspace{2.0mm}

\noindent
\begin{tabularx}{\linewidth}{|X|}
\hline
\vspace{0.5mm}
\textbf{Câu 13 (2,5 điểm).} Cho tam giác $ABC$ vuông tại $A$, đường cao $AH$ ($H \in BC$).\par
\textbf{a) (1,5 điểm)} Cho $AB = 9\text{ cm}$ và $AC = 12\text{ cm}$. Tính độ dài cạnh huyền $BC$, đường cao $AH$ và số đo góc $B$ (làm tròn đến độ).\par
\textbf{b) (1,0 điểm)} Kẻ $HD \perp AB$ tại $D$ và $HE \perp AC$ tại $E$. Chứng minh $\Delta ADE \backsim \Delta ACB$.\par
\vspace{0.4mm}
\noindent\textit{Bài giải:}\dongke{18}
\\ \hline
\end{tabularx}

\newpage
% =============================================================================
% ========================= TOÁN 9 - TUẦN 08 - BUỔI 02: MÃ ĐỀ 102 =============
% ========================= MẶT TRƯỚC (TRANG 1/2) ==============================
% =============================================================================

\noindent
\begin{tabularx}{\linewidth}{|X|c|}
\hline
\textbf{PHIẾU KIỂM TRA ĐỊNH KỲ TOÁN 9 - TUẦN 08 - BUỔI 02} & \textbf{Mã đề: 102 - Trang 1/2} \\
\hline
\end{tabularx}

\vspace{1.0mm}

\noindent
\begin{tabularx}{\linewidth}{|p{3.2cm}|X|}
\hline
\centering\textbf{ĐIỂM SỐ} & \centering\textbf{LỜI NHẬN XÉT CỦA THẦY CÔ} \tabularnewline
\hline
\rule{0pt}{1.8cm} & \tabularnewline
\hline
\end{tabularx}

\vspace{1.2mm}
\noindent\textbf{Câu 1 (0,25 điểm).} Giá trị của biểu thức $\sqrt{64} - \sqrt{25}$ bằng: \shortans{3}\par\vspace{0.5mm}
\noindent\textbf{Câu 2 (0,25 điểm).} Biểu thức $\sqrt{3x - 12}$ xác định khi: \shortans{x \ge 4}\par\vspace{0.5mm}
\noindent\textbf{Câu 3 (0,25 điểm).} Nghiệm của bất phương trình $2x - 8 \ge 0$ là: \shortans{x \ge 4}\par\vspace{0.5mm}
\noindent\textbf{Câu 4 (0,25 điểm).} Nghiệm của hệ $\begin{cases} 2x - y = 4 \\ x + y = 5 \end{cases}$ là $(x; y) =$ \shortans{(3; 2)}\par\vspace{0.5mm}
\noindent\textbf{Câu 5 (0,25 điểm).} Cho $\Delta ABC$ vuông tại $A$ ($AB=3, AC=4$). Khi đó $\cos C =$ \shortans{\dfrac{4}{5}}\par\vspace{0.5mm}
\noindent\textbf{Câu 6 (0,25 điểm).} $\Delta ABC$ vuông tại $A$ ($BC = 8, \widehat{B}=60^\circ$). Cạnh $AC =$ \shortans{4\sqrt{3}\text{ cm}}\par\vspace{0.5mm}
\noindent\textbf{Câu 7 (0,25 điểm).} Rút gọn biểu thức $\sqrt{(2 - \sqrt{5})^2}$ được: \shortans{\sqrt{5} - 2}\par\vspace{0.5mm}
\noindent\textbf{Câu 8 (0,25 điểm).} Giá trị của biểu thức $\cos^2 25^\circ + \cos^2 65^\circ$ bằng: \shortans{1}\par\vspace{0.5mm}
\noindent\textbf{Câu 9 (0,25 điểm).} Cột đèn cao $6\text{ m}$ đổ bóng dài $6\text{ m}$. Góc của tia nắng là: \shortans{45^\circ}\par\vspace{0.5mm}
\noindent\textbf{Câu 10 (0,25 điểm).} Nghiệm của phương trình $\sqrt{2x + 5} = 5$ là $x =$ \shortans{10}\par\vspace{0.5mm}
\vspace{0.8mm}
\noindent\textbf{Câu 11 (2,5 điểm). Đại số tổng hợp:}
\vspace{0.6mm}

\noindent
\begin{tabularx}{\linewidth}{|X|X|}
\hline
\oBaiTap[22]{a) (1,25 điểm)}{Giải hệ phương trình:\par $\begin{cases} 3x - y = 7 \\ 2x + y = 8 \end{cases}$} & 
\oBaiTap[22]{b) (1,25 điểm)}{Rút gọn với $x > 0, x \ne 4$:\par $B = \left(\dfrac{\sqrt{x}}{\sqrt{x}-2} - \dfrac{4}{x-2\sqrt{x}}\right) : \dfrac{\sqrt{x}+2}{x-4}$} \\
\hline
\end{tabularx}

\newpage
% =============================================================================
% ========================= TOÁN 9 - TUẦN 08 - BUỔI 02: MÃ ĐỀ 102 =============
% ========================= MẶT SAU (TRANG 2/2) ================================
% =============================================================================

\noindent
\begin{tabularx}{\linewidth}{|X|c|}
\hline
\textbf{TOÁN 9 - TUẦN 08 - BUỔI 02} & \textbf{Mã đề: 102 - Trang 2/2} \\
\hline
\end{tabularx}

\vspace{1.2mm}

\noindent
\begin{tabularx}{\linewidth}{|X|}
\hline
\vspace{0.5mm}
\textbf{Câu 12 (2,5 điểm).} Một đội xe theo kế hoạch phải chở $60$ tấn hàng để phục vụ xây dựng trường học. Khi thực hiện, có $2$ xe được điều đi làm nhiệm vụ khác nên mỗi xe còn lại phải chở thêm $1$ tấn hàng nữa mới hết số hàng. Hỏi ban đầu đội xe có bao nhiêu chiếc xe? (Biết mỗi xe chở khối lượng hàng như nhau).\par
\vspace{0.4mm}
\noindent\textit{Bài giải:}\dongke{16}
\\ \hline
\end{tabularx}

\vspace{2.0mm}

\noindent
\begin{tabularx}{\linewidth}{|X|}
\hline
\vspace{0.5mm}
\textbf{Câu 13 (2,5 điểm).} Cho tam giác $ABC$ vuông tại $A$, đường cao $AH$ ($H \in BC$).\par
\textbf{a) (1,5 điểm)} Cho $AB = 12\text{ cm}$ và $AC = 16\text{ cm}$. Tính độ dài cạnh huyền $BC$, đường cao $AH$ và số đo góc $B$ (làm tròn đến độ).\par
\textbf{b) (1,0 điểm)} Kẻ $HD \perp AB$ tại $D$ và $HE \perp AC$ tại $E$. Chứng minh $\Delta ADE \backsim \Delta ACB$.\par
\vspace{0.4mm}
\noindent\textit{Bài giải:}\dongke{18}
\\ \hline
\end{tabularx}

\newpage
% =============================================================================
% ========================= TOÁN 9 - TUẦN 08 - BUỔI 02: MÃ ĐỀ 103 =============
% ========================= MẶT TRƯỚC (TRANG 1/2) ==============================
% =============================================================================

\noindent
\begin{tabularx}{\linewidth}{|X|c|}
\hline
\textbf{PHIẾU KIỂM TRA ĐỊNH KỲ TOÁN 9 - TUẦN 08 - BUỔI 02} & \textbf{Mã đề: 103 - Trang 1/2} \\
\hline
\end{tabularx}

\vspace{1.0mm}

\noindent
\begin{tabularx}{\linewidth}{|p{3.2cm}|X|}
\hline
\centering\textbf{ĐIỂM SỐ} & \centering\textbf{LỜI NHẬN XÉT CỦA THẦY CÔ} \tabularnewline
\hline
\rule{0pt}{1.8cm} & \tabularnewline
\hline
\end{tabularx}

\vspace{1.2mm}
\noindent\textbf{Câu 1 (0,25 điểm).} Giá trị của biểu thức $\sqrt{81} - \sqrt{36}$ bằng: \shortans{3}\par\vspace{0.5mm}
\noindent\textbf{Câu 2 (0,25 điểm).} Biểu thức $\sqrt{4x - 8}$ xác định khi: \shortans{x \ge 2}\par\vspace{0.5mm}
\noindent\textbf{Câu 3 (0,25 điểm).} Nghiệm của bất phương trình $4x - 12 \ge 0$ là: \shortans{x \ge 3}\par\vspace{0.5mm}
\noindent\textbf{Câu 4 (0,25 điểm).} Nghiệm của hệ $\begin{cases} 3x - y = 5 \\ x + y = 3 \end{cases}$ là $(x; y) =$ \shortans{(2; 1)}\par\vspace{0.5mm}
\noindent\textbf{Câu 5 (0,25 điểm).} Cho $\Delta ABC$ vuông tại $A$ ($AB=5, AC=12$). Khi đó $\sin C =$ \shortans{\dfrac{5}{13}}\par\vspace{0.5mm}
\noindent\textbf{Câu 6 (0,25 điểm).} $\Delta ABC$ vuông tại $A$ ($BC = 6, \widehat{B}=60^\circ$). Cạnh $AC =$ \shortans{3\sqrt{3}\text{ cm}}\par\vspace{0.5mm}
\noindent\textbf{Câu 7 (0,25 điểm).} Rút gọn biểu thức $\sqrt{(1 - \sqrt{2})^2}$ được: \shortans{\sqrt{2} - 1}\par\vspace{0.5mm}
\noindent\textbf{Câu 8 (0,25 điểm).} Giá trị của biểu thức $\sin^2 18^\circ + \sin^2 72^\circ$ bằng: \shortans{1}\par\vspace{0.5mm}
\noindent\textbf{Câu 9 (0,25 điểm).} Tháp cao $50\text{ m}$ đổ bóng dài $50\text{ m}$. Góc của tia nắng là: \shortans{45^\circ}\par\vspace{0.5mm}
\noindent\textbf{Câu 10 (0,25 điểm).} Nghiệm của phương trình $\sqrt{4x + 9} = 5$ là $x =$ \shortans{4}\par\vspace{0.5mm}
\vspace{0.8mm}
\noindent\textbf{Câu 11 (2,5 điểm). Đại số tổng hợp:}
\vspace{0.6mm}

\noindent
\begin{tabularx}{\linewidth}{|X|X|}
\hline
\oBaiTap[22]{a) (1,25 điểm)}{Giải hệ phương trình:\par $\begin{cases} 2x + y = 11 \\ 3x - y = 4 \end{cases}$} & 
\oBaiTap[22]{b) (1,25 điểm)}{Rút gọn với $x > 0, x \ne 9$:\par $B = \left(\dfrac{\sqrt{x}}{\sqrt{x}-3} - \dfrac{9}{x-3\sqrt{x}}\right) : \dfrac{\sqrt{x}+3}{x-9}$} \\
\hline
\end{tabularx}

\newpage
% =============================================================================
% ========================= TOÁN 9 - TUẦN 08 - BUỔI 02: MÃ ĐỀ 103 =============
% ========================= MẶT SAU (TRANG 2/2) ================================
% =============================================================================

\noindent
\begin{tabularx}{\linewidth}{|X|c|}
\hline
\textbf{TOÁN 9 - TUẦN 08 - BUỔI 02} & \textbf{Mã đề: 103 - Trang 2/2} \\
\hline
\end{tabularx}

\vspace{1.2mm}

\noindent
\begin{tabularx}{\linewidth}{|X|}
\hline
\vspace{0.5mm}
\textbf{Câu 12 (2,5 điểm).} Một đội xe dự định chở $80$ tấn hàng cứu trợ đồng bào bị ảnh hưởng do bão. Lúc sắp khởi hành, có $4$ xe được điều đi làm việc khác nên mỗi xe còn lại phải chở nhiều hơn dự định $1$ tấn hàng mới hết số hàng. Hỏi ban đầu đội xe có bao nhiêu chiếc xe? (Biết mỗi xe chở khối lượng hàng như nhau).\par
\vspace{0.4mm}
\noindent\textit{Bài giải:}\dongke{16}
\\ \hline
\end{tabularx}

\vspace{2.0mm}

\noindent
\begin{tabularx}{\linewidth}{|X|}
\hline
\vspace{0.5mm}
\textbf{Câu 13 (2,5 điểm).} Cho tam giác $ABC$ vuông tại $A$, đường cao $AH$ ($H \in BC$).\par
\textbf{a) (1,5 điểm)} Cho $AB = 6\text{ cm}$ và $AC = 8\text{ cm}$. Tính độ dài cạnh huyền $BC$, đường cao $AH$ và số đo góc $B$ (làm tròn đến độ).\par
\textbf{b) (1,0 điểm)} Kẻ $HD \perp AB$ tại $D$ và $HE \perp AC$ tại $E$. Chứng minh $\Delta ADE \backsim \Delta ACB$.\par
\vspace{0.4mm}
\noindent\textit{Bài giải:}\dongke{18}
\\ \hline
\end{tabularx}

\newpage
% =============================================================================
% ========================= TOÁN 9 - TUẦN 08 - BUỔI 02: MÃ ĐỀ 104 =============
% ========================= MẶT TRƯỚC (TRANG 1/2) ==============================
% =============================================================================

\noindent
\begin{tabularx}{\linewidth}{|X|c|}
\hline
\textbf{PHIẾU KIỂM TRA ĐỊNH KỲ TOÁN 9 - TUẦN 08 - BUỔI 02} & \textbf{Mã đề: 104 - Trang 1/2} \\
\hline
\end{tabularx}

\vspace{1.0mm}

\noindent
\begin{tabularx}{\linewidth}{|p{3.2cm}|X|}
\hline
\centering\textbf{ĐIỂM SỐ} & \centering\textbf{LỜI NHẬN XÉT CỦA THẦY CÔ} \tabularnewline
\hline
\rule{0pt}{1.8cm} & \tabularnewline
\hline
\end{tabularx}

\vspace{1.2mm}
\noindent\textbf{Câu 1 (0,25 điểm).} Giá trị của biểu thức $\sqrt{100} - \sqrt{49}$ bằng: \shortans{3}\par\vspace{0.5mm}
\noindent\textbf{Câu 2 (0,25 điểm).} Biểu thức $\sqrt{5x - 25}$ xác định khi: \shortans{x \ge 5}\par\vspace{0.5mm}
\noindent\textbf{Câu 3 (0,25 điểm).} Nghiệm của bất phương trình $5x - 15 \ge 0$ là: \shortans{x \ge 3}\par\vspace{0.5mm}
\noindent\textbf{Câu 4 (0,25 điểm).} Nghiệm của hệ $\begin{cases} 2x - y = 5 \\ x + y = 4 \end{cases}$ là $(x; y) =$ \shortans{(3; 1)}\par\vspace{0.5mm}
\noindent\textbf{Câu 5 (0,25 điểm).} Cho $\Delta ABC$ vuông tại $A$ ($AB=5, AC=12$). Khi đó $\cos C =$ \shortans{\dfrac{12}{13}}\par\vspace{0.5mm}
\noindent\textbf{Câu 6 (0,25 điểm).} $\Delta ABC$ vuông tại $A$ ($BC = 12, \widehat{B}=60^\circ$). Cạnh $AC =$ \shortans{6\sqrt{3}\text{ cm}}\par\vspace{0.5mm}
\noindent\textbf{Câu 7 (0,25 điểm).} Rút gọn biểu thức $\sqrt{(3 - \sqrt{10})^2}$ được: \shortans{\sqrt{10} - 3}\par\vspace{0.5mm}
\noindent\textbf{Câu 8 (0,25 điểm).} Giá trị của biểu thức $\cos^2 40^\circ + \cos^2 50^\circ$ bằng: \shortans{1}\par\vspace{0.5mm}
\noindent\textbf{Câu 9 (0,25 điểm).} Ống khói cao $25\text{ m}$ đổ bóng dài $25\text{ m}$. Góc của tia nắng là: \shortans{45^\circ}\par\vspace{0.5mm}
\noindent\textbf{Câu 10 (0,25 điểm).} Nghiệm của phương trình $\sqrt{5x + 4} = 7$ là $x =$ \shortans{9}\par\vspace{0.5mm}
\vspace{0.8mm}
\noindent\textbf{Câu 11 (2,5 điểm). Đại số tổng hợp:}
\vspace{0.6mm}

\noindent
\begin{tabularx}{\linewidth}{|X|X|}
\hline
\oBaiTap[22]{a) (1,25 điểm)}{Giải hệ phương trình:\par $\begin{cases} 4x - y = 9 \\ 2x + y = 9 \end{cases}$} & 
\oBaiTap[22]{b) (1,25 điểm)}{Rút gọn với $x > 0, x \ne 16$:\par $B = \left(\dfrac{\sqrt{x}}{\sqrt{x}-4} - \dfrac{16}{x-4\sqrt{x}}\right) : \dfrac{\sqrt{x}+4}{x-16}$} \\
\hline
\end{tabularx}

\newpage
% =============================================================================
% ========================= TOÁN 9 - TUẦN 08 - BUỔI 02: MÃ ĐỀ 104 =============
% ========================= MẶT SAU (TRANG 2/2) ================================
% =============================================================================

\noindent
\begin{tabularx}{\linewidth}{|X|c|}
\hline
\textbf{TOÁN 9 - TUẦN 08 - BUỔI 02} & \textbf{Mã đề: 104 - Trang 2/2} \\
\hline
\end{tabularx}

\vspace{1.2mm}

\noindent
\begin{tabularx}{\linewidth}{|X|}
\hline
\vspace{0.5mm}
\textbf{Câu 12 (2,5 điểm).} Một đội xe theo kế hoạch phải chở $90$ tấn nông sản đi tiêu thụ tại các siêu thị. Khi thực hiện, do có $3$ xe được điều chuyển đi nơi khác nên mỗi xe còn lại phải chở thêm $1$ tấn hàng mới hết số hàng. Hỏi ban đầu đội xe có bao nhiêu chiếc xe? (Biết mỗi xe chở khối lượng hàng như nhau).\par
\vspace{0.4mm}
\noindent\textit{Bài giải:}\dongke{16}
\\ \hline
\end{tabularx}

\vspace{2.0mm}

\noindent
\begin{tabularx}{\linewidth}{|X|}
\hline
\vspace{0.5mm}
\textbf{Câu 13 (2,5 điểm).} Cho tam giác $ABC$ vuông tại $A$, đường cao $AH$ ($H \in BC$).\par
\textbf{a) (1,5 điểm)} Cho $AB = 15\text{ cm}$ và $AC = 20\text{ cm}$. Tính độ dài cạnh huyền $BC$, đường cao $AH$ và số đo góc $B$ (làm tròn đến độ).\par
\textbf{b) (1,0 điểm)} Kẻ $HD \perp AB$ tại $D$ và $HE \perp AC$ tại $E$. Chứng minh $\Delta ADE \backsim \Delta ACB$.\par
\vspace{0.4mm}
\noindent\textit{Bài giải:}\dongke{18}
\\ \hline
\end{tabularx}

\end{document}
